(the sum extends over every particle in the system; for brevity, the label identifying a particle is suppressed). The matrix

\SourcePage{343}{346}
of this operator---as with any polar vector (see \S~30)---can have nonzero entries only for transitions between states of opposite parity. Consequently, all its diagonal entries are necessarily zero. Equivalently, in a stationary state the mean dipole moment of any particle system (an atom, for example) vanishes\footnote{To prevent a possible misunderstanding, we stress that this statement concerns either a closed particle system or one placed in an external electric field having a center of symmetry. Thus, if the nuclei are regarded as ``fixed,'' the general statement just made applies to the electron system of an atom, but not to that of a molecule.

It is also assumed that the energy level has no extra (``accidental'') degeneracy beyond the degeneracy associated with the orientations of the total angular momentum. Otherwise one can form wave functions for stationary states that have no definite parity, and the corresponding diagonal entries of the dipole moment need not vanish.}.
